\section{Conclusion}
\label{sec:conclusion}

We introduce \textsc{TIDE-Bench}, a benchmark for evaluating how text-to-SQL agents handle two underexamined  patterns in conversational interaction: chain ambiguity and user intent drift. Built around a paired factorial design over 514 anchor SQLs from BIRD, \textsc{TIDE-Bench} contains 1{,}542 samples constructed through a four-stage automated pipeline with multi-stage validation and human verification. Our experiments on 12 LLMs reveal three failure modes. Chain ambiguity exposes an identification bottleneck unaffected by ask behavior, while user intent drift reveals a wide gap between recognition and resolution, and the two failures overlap when they co-occur. The paired design attributes these gaps to interaction structure rather than uncontrolled differences in schemas or query difficulty. These findings suggest that future work on conversational text-to-SQL should look beyond surface clarification behaviors toward modeling identification capacity and tracking dialogue-level commitments. This requires agents to surface dependent ambiguities and revise prior commitments without inheriting earlier errors. We hope \textsc{TIDE-Bench} provides a diagnostic tool for both research and development in conversational text-to-SQL.